\section{Numerical implementation}\label{sec:numerics}
\subsection{Scales, penalty and conditional least squares}\label{sec:fit-details}
The scale of transfer $(o,i)$ in Eq.~\eqref{eq:objective} is
\[
 d_{oi}=\max\Bigl\{\rho_{oi},\ 10^{-3}\operatorname{median}_{o,i}\rho_{oi},\
 10^{-12}\ \mathrm{m\,s^{-2}/N}\Bigr\},\qquad
 \rho_{oi}=\Bigl(N^{-1}\sum_{\ell}|\widehat H_{\ell,oi}|^{2}\Bigr)^{1/2},
\]
computed once on the training record and frozen for every later record.
For the penalty, let $E_{\rm g}$ be the geometric extension of the $h$ port
directions over the displayed mesh, built by minimum-roughness interpolation
and corrected to reproduce five rigid patterns in the laboratory, and
let $Q_{\rm g}\succeq0$ be the surface-gradient roughness operator assembled
from the mesh connectivity; the chain uses first differences along the chain.
Then $G=E_{\rm g}^{\top}Q_{\rm g}E_{\rm g}/\|E_{\rm g}^{\top}Q_{\rm g}E_{\rm g}\|_2$
and $\beta=\alpha/(r\chi_0^{2})$ with $\chi_0^{2}=\operatorname{median}_{o,i}d_{oi}$
and $\alpha=10^{-5}$ in both examples. A zero roughness gives a zero penalty.
Here $\|\cdot\|_2$ is the Euclidean vector norm or the spectral matrix norm
and $\|\cdot\|_F$ the Frobenius norm.
For disjoint ports the scaling
$(R_{\rm o}W,R_{\rm i}W)\mapsto(c^{-1}R_{\rm o}W,cR_{\rm i}W)$ with $c>0$
leaves $H_\theta$ unchanged. The penalty balances this reciprocal scale.

For disjoint ports and fixed $(R_{\rm i}W,\omega,\Gamma)$, $H_\theta(s_\ell)=R_{\rm o}WD_\ell$ with
$D_\ell=s_\ell^{2}\Delta_\theta(s_\ell)^{-1}W^{\top}R_{\rm i}^{\top}$, so the response
rows minimize a real linear least-squares objective: the real and imaginary
parts of $(R_{\rm o}WD_\ell-\widehat H_\ell)_{oi}/d_{oi}$, stacked over $\ell$ and $i$,
plus the quadratic penalty, whose blocks couple response rows with each other
and with the fixed force rows. Writing $G=\Xi^{\top}\Xi$ turns the penalty
into augmented residual rows; QR compression of the frequency equations and a
right-preconditioned LSMR solve then avoid assembling the coupled normal matrix.
Variable projection \citep{golub2003,oleary2013} thus eliminates the response
rows.

The nonlinear derivatives use $d(\Delta^{-1})=-\Delta^{-1}(d\Delta)\Delta^{-1}$
and the envelope property of the conditional minimum \citep{golub2003,oleary2013}.
Positivity is built into the parameterization,
\begin{equation}
 \omega_j=\omega_{\rm ref}\exp(\vartheta_j),\qquad
 \Gamma=\omega_{\rm ref}(\Upsilon\Upsilon^{\top}+10^{-7}I_r),\qquad
 \Upsilon_{jj}=\exp(\gamma_{jj}),\quad
 \Upsilon_{jk}=\gamma_{jk}\ (0<j-k\le b),
 \label{eq:positive}
\end{equation}
with $\omega_{\rm ref}=2\pi\max_\ell\nu_\ell$, all other entries of the lower
triangular $\Upsilon$ zero, and $b$ the bandwidth in optimizer order. The
real variables $\vartheta_j,\gamma_{jk}$ and the force rows $R_{\rm i}W$ are
updated by L-BFGS-B \citep{byrd1995} within
$2\pi\max(0.05\min_\ell\nu_\ell,0.001\ \mathrm{Hz})\le\omega_j\le3\omega_{\rm ref}$,
factor diagonals in $[10^{-5},2]$ and off-diagonals in $[-2,2]$. The chain
uses $b=2$ (a full factor) and the laboratory $b=24$. Exported frequency
orderings permute $W,\Omega^{2},\Gamma$ together and are undone before
completion. The laboratory continuation started from 13.40\% training and
16.10\% development error. Directional reduced-gradient checks at two
laboratory states gave relative differences $3.17\times10^{-7}$ and
$1.24\times10^{-8}$. The order $r$ and the bandwidth $b$ are fixed before
fitting. The pole expansion in Eq.~\eqref{eq:poleresidue} makes each
evaluation linear in $N$. Hereafter $\theta$ denotes the returned fitted
parameters. No global optimality is claimed, and the examples report the
stopping diagnostics.

\subsection{Selecting the hidden fields with sparse solves}\label{sec:field-selection}
The laboratory trial space $\mathcal U$ has two parts. The first holds the
static deflections $X=K_0^{-1}A^{\top}$ under unit port forces. Their port
values form the port flexibility $AX=AK_0^{-1}A^{\top}$, which is positive
definite because $A$ has full row rank. Recombining the deflections to unit
port values gives the static lift
\begin{equation}
 E=K_0^{-1}A^{\top}(AK_0^{-1}A^{\top})^{-1},\qquad AE=I_h.
 \label{eq:lift}
\end{equation}
Its column $Ew$ is the field of least strain energy with port values $w$.
Since $K_0E=A^{\top}(AX)^{-1}$, every $v\in\ker A$ has
$E^{\top}K_0v=(AX)^{-1}Av=0$. Every $t$ with $At=w$ is $t=Ew+v$ with
$v\in\ker A$, and then $\|t\|_{K_0}^2=\|Ew\|_{K_0}^2+\|v\|_{K_0}^2$.
The second part holds $d$ port-invisible fields $N$, selected below and
normalized by
\begin{equation}
 AN=0,\qquad N^{\top}M_0N=I_d.
 \label{eq:hidden}
\end{equation}
Every candidate map in $\mathcal U$ then reads
\begin{equation}
 T=EW+NZ,\qquad Z\in\R^{d\times r},
 \label{eq:completion}
\end{equation}
and satisfies $AT=W$ for every $Z$. This form is only a computational
device. Theorem~\ref{thm:completion} does not depend on it. A change of
basis of the fields with the inverse change of $Z$ leaves $T$ unchanged.
The static map
$EW$ also sets the margin $\delta=10^{-3}\sigma_1(REW)$.

Here $R\in\R^{n_{\rm d}\times n}$ selects the $n_{\rm d}=7988$ displayed
translations. Let $M_{\rm d}$ be the positive diagonal of $M_0$ on those coordinates, and set
$D_{\rm d}=M_{\rm d}^{1/2}$. The laboratory prior assumes static equilibrium in
the omitted coordinates. Among fields satisfying that assumption and
$Au=0$, the retained directions have the smallest values of
\[
 \frac{u^{\top}K_0u}{(Ru)^{\top}M_{\rm d}(Ru)}.
\]
This selects mechanically compliant invisible fields. A generic null-space
SVD enforces $Au=0$, but supplies no such stiffness ordering.

The selection requires neither a dense null-space basis nor a condensed
stiffness. Factor $K_0$ once and compute
\[
 K_0X=A^{\top},\qquad G_A=AX,\qquad EG_A=X.
\]
The last equation is a solve with the $h\times h$ constraint matrix.
The projection $\mathcal Pv=v-E(Av)$ removes registered motion. It is applied
without forming $\mathcal P=I-EA$. Define the symmetric positive semidefinite
selection operator by
\begin{equation}
 \boxed{\quad K_0w=R^{\top}D_{\rm d}v,\qquad
 \mathcal Hv=D_{\rm d}R\bigl(w-E(Aw)\bigr).\quad}
 \label{eq:selection-action}
\end{equation}
Its largest positive eigenpairs $\mathcal Hy_k=\mu_ky_k$ correspond to the
smallest stiffness-selection eigenvalues $\varpi_k=1/\mu_k$.
These are eigenvalues of the statically relaxed, displayed-metric problem.
Their square roots are not the full FE model's natural frequencies.
Appendix~\ref{sec:field-details} derives this equivalence.

The same solves recover each full field:
\begin{equation}
 \widetilde n_k=\frac1{\mu_k}
 (K_0^{-1}-EAK_0^{-1})R^{\top}D_{\rm d}y_k,
 \qquad A\widetilde n_k=0,\quad
 R\widetilde n_k=D_{\rm d}^{-1}y_k.
 \label{eq:field-recovery}
\end{equation}
Because $A$ has support only on displayed translations, the right-hand
side of $K_0\widetilde n_k$ vanishes on omitted coordinates. Thus the
recovered field satisfies their static equilibrium, without a second
factorization. Keep complete clusters at the retention boundary.
For the resulting $\widetilde N$, only the small retained Gram matrix is
normalized:
\begin{equation}
 G_N=\widetilde N^{\top}M_0\widetilde N=LL^{\top},\qquad
 N=\widetilde NL^{-\top}.
 \label{eq:retained-normalization}
\end{equation}
Here $L$ is lower triangular. Selection uses $M_{\rm d}$. Normalization uses
the full $M_0$, including the recovered coordinates. Substituting a condensed
mass or selecting full-mass modes would choose a different approximation
space. Triangular normalization and a subsequent orthogonal rotation only
change coordinates within the selected space.

The completion runs as follows.
\begin{enumerate}
\item Supply $K_0,M_0,A,R,M_{\rm d}$, the fitted $(W,\omega,\Gamma)$,
the retained-field and boundary-cluster rule and the independence margin.
\item Factor $K_0$ once. Solve for $X=K_0^{-1}A^{\top}$ and factor
$G_A=AX$ to construct $E$. Verify $AE=I_h$.
\item Apply $\mathcal H$ by Eq.~\eqref{eq:selection-action}, using the same
stiffness factor and the small constraint correction.
\item Compute only the required largest positive eigenpairs, with extra
pairs to detect a cutoff cluster in $\varpi=1/\mu$.
\item Recover the full $\widetilde N$ by Eq.~\eqref{eq:field-recovery} in
bounded batches. Check its ports, displayed fields and omitted equilibrium.
\item Factor the retained mass Gram to obtain $N$. Verify
$AN=0$ and $N^{\top}M_0N=I_d$.
\item Project $K_0$, $M_0$ and $M_0K_0^{-1}M_0$ onto $[E,N]$, solve the
equality-only problem of Appendix~\ref{sec:constrained} and check $\sigma_r(RT)$. If the margin fails,
run the constrained iteration and report its feasibility and stopping
residuals. In either case verify $AT=W$.
\end{enumerate}

\subsection{Static condensation and the hidden fields}\label{sec:field-details}
In the laboratory the hidden fields are built on the displayed translations
and condensed over the rest. Write $u_{\rm d}=Ru$ for the displayed
translations and let $u_{\rm o}$ collect the omitted translations and
rotations. In the order $(u_{\rm d},u_{\rm o})$ write the blocks of $K_0$ as
$K_{\rm dd},K_{\rm do},K_{\rm od},K_{\rm oo}$ with $K_{\rm do}=K_{\rm od}^{\top}$.
Completing the block square,
\begin{equation}
 u^{\top}K_0u=u_{\rm d}^{\top}
 (K_{\rm dd}-K_{\rm do}K_{\rm oo}^{-1}K_{\rm od})u_{\rm d}
 +(u_{\rm o}+K_{\rm oo}^{-1}K_{\rm od}u_{\rm d})^{\top}K_{\rm oo}
 (u_{\rm o}+K_{\rm oo}^{-1}K_{\rm od}u_{\rm d}),
 \label{eq:block-square}
\end{equation}
so for given $u_{\rm d}$ the energy is least at
$u_{\rm o}=-K_{\rm oo}^{-1}K_{\rm od}u_{\rm d}$. The corresponding static lift
\begin{equation}
 u=\Lambda u_{\rm d},\qquad R\Lambda=I,\qquad
 \Lambda=\begin{bmatrix}I\\-K_{\rm oo}^{-1}K_{\rm od}\end{bmatrix}_{(\rm d,o)}
 \label{eq:static-lift}
\end{equation}
is a right inverse of $R$: the static expansion underlying Guyan reduction
\citep{guyan1965}. Its image minimizes elastic energy at fixed displayed
translations.
The condensed stiffness is $\Lambda^{\top}K_0\Lambda$, the Schur complement in
Eq.~\eqref{eq:block-square}, and the registration on displayed coordinates is
$A\Lambda$. Since the ports involve only displayed translations, the lift $E$
of Eq.~\eqref{eq:lift} also lies in $\range\Lambda$.

The selection metric $M_{\rm d}=\diag((M_0)_{\rm dd})$ is neither the
condensed mass $\Lambda^{\top}M_0\Lambda$ nor a full modal problem. Let
$U_A$ be any basis of $\ker(A\Lambda)$ with independent columns. An
anchor-invisible displayed field is then $U_Av$. With
$K_N=U_A^{\top}\Lambda^{\top}K_0\Lambda U_A$ and $M_N=U_A^{\top}M_{\rm d}U_A$,
both positive definite, the field problem is the symmetric generalized
eigenproblem
\begin{equation}
 K_Nv_k=\varpi_kM_Nv_k,\qquad v_i^{\top}M_Nv_j=\delta_{ij},\qquad
 \widetilde n_k=\Lambda U_Av_k,
 \label{eq:fe-fields}
\end{equation}
with mass-orthonormal reduced coordinates $v_k$ in this explicit
representation. The smallest selection
eigenvalues $\varpi_k$ minimize stiffness relative to $M_{\rm d}$. Their
square roots need not equal natural frequencies of the full FE model. The
implementation never forms $U_A$ or a condensed matrix. The constrained
inverse $U_AK_N^{-1}U_A^{\top}=R(K_0^{-1}-EAK_0^{-1})R^{\top}$ is applied
through sparse solves with $K_0$, scaled on both sides by $M_{\rm d}^{1/2}$,
and its largest eigenvalues $1/\varpi_k$ are computed by a Lanczos iteration.
The Euclidean-orthonormal eigenvectors of Eq.~\eqref{eq:selection-action}
are $y_k=M_{\rm d}^{1/2}U_Av_k$. Both representations describe the same fields.

To see the inverse identity directly, for a displayed load $f$ set
$u=(K_0^{-1}-EAK_0^{-1})R^{\top}f$. Then $Au=0$ and
$K_0u=R^{\top}f-A^{\top}G_A^{-1}AK_0^{-1}R^{\top}f$.
The omitted rows have zero right-hand side, hence $u=\Lambda Ru$.
Testing against $\Lambda U_A$ gives $K_Nv=U_A^{\top}f$ for $Ru=U_Av$,
which proves the displayed constrained-inverse formula. In particular,
$K_0^{-1}-EAK_0^{-1}$ is symmetric even though $I-EA$ is generally not
Euclidean-orthogonal. It is the $K_0$-orthogonal projection followed by
the elastic solve.

Initially $d=\min(r,\dim\ker(A\Lambda))$ fields are retained, with complete
boundary clusters added when adjacent $\varpi_k$ differ by at most
$10^{-8}\max(1,|\varpi_k|)$ in SI units. The selected fields are recovered
by Eq.~\eqref{eq:field-recovery}, which applies $\Lambda$ implicitly with
the full stiffness factor, and projected by $I_n-EA$ to remove roundoff
anchor components. The retained-Gram Cholesky factor in
Eq.~\eqref{eq:retained-normalization} then normalizes them in the full mass.
A Rayleigh--Ritz rotation
diagonalizes $N^{\top}K_0N$ within the fixed span and the largest-magnitude
entry fixes each column's sign. This preserves the selected span and gives
Eq.~\eqref{eq:hidden}. It does not turn the selection into a full-$M_0$
eigenproblem. The chain, with
$n<100$, uses the dense null space of $A$ directly.

The sparse implementation uses equilibrated SuperLU factors and two
residual-refinement steps per elastic solve. The partial symmetric
eigensolver uses tolerance $10^{-10}$, the deterministic start
$v_i=\sin(i\sqrt2)$, and initially requests eight extra eigenpairs.
The request grows if the complete boundary cluster has not yet been
bracketed. A nonconverged eigensolve is rejected explicitly and never
returns a partial dictionary. Recovery and projected dual-energy
solves use at most 32 right-hand sides per batch. The retained Gram is
checked for numerical rank before normalization. A full-rank spectral
factor is available if Cholesky fails. No ridge or field deletion hides rank
loss. The final Ritz rotation is retained for its small
cost, stiffness ordering and diagonal hidden stiffness block. It does not
diagonalize the mass and dual-energy cross blocks with $E$.

The stiffness factor depends only on $K_0$. A change in $A$ recomputes its
small constraint factor and $E$. A change in $R,M_{\rm d},d$ or the cutoff
rule invalidates field selection. Full $M_0$ also determines normalization
and projected operators. Fitted $W,\omega,\Gamma$ do not invalidate the
dictionary unless the specified field count changes. Cached files contain
content signatures of the nominal matrices, registration, displayed rows
and selection rules. In-memory contexts own their inputs and reject a
changed stiffness instead of reusing a stale factor.

\paragraph{What an SVD alternative would change.}
Write $A_{\rm d}=AR^{\top}$. If its nonzero-column support has size $s$, a
column permutation $\Pi_s$ gives $A_{\rm d}\Pi_s=[A_s\;0]$. A small SVD can
compute $V_s=\operatorname{null}(A_s)$ and represent the displayed null
space as
\[
 U_A=\Pi_s\begin{bmatrix}V_s&0\\0&I_{n_{\rm d}-s}\end{bmatrix}.
\]
This representation need never be materialized. The identity block acts
directly on unsupported coordinates. It enforces the same interpolated
ports but still needs the mechanical selection problem.
Alternatively, a full-column-rank nominal candidate matrix
$\Phi\in\R^{n\times n_\Phi}$ gives $N=\Phi\operatorname{null}(A\Phi)$ and
\[
 \range N=\range\Phi\cap\ker A,\qquad
 d=n_\Phi-\rank(A\Phi).
\]
Thus $n_\Phi=417$ yields $d=384$ only if $\rank(A\Phi)=33$. This is a
different approximation space. The dimension count alone neither
demonstrates spatial accuracy nor preserves the static assumption. This
paper uses neither candidate modes nor response-based enrichment.

\subsection{Equality-only least squares and constrained refinement}\label{sec:constrained}
In the basis $[E,N]$ only three small matrices are assembled:
$\widehat K=[E,N]^{\top}K_0[E,N]$, $\widehat M=[E,N]^{\top}M_0[E,N]$ and
$\widehat H=[E,N]^{\top}M_0K_0^{-1}M_0[E,N]$, the last with the reused
sparse stiffness factor. For $T=[E,N]X$ with $X=[W;Z]$,
\[
 \tr(G^{\top}K_0^{-1}G\Phi_q)=\tr(X^{\top}\widehat KX\Phi_q)
 -\tr\bigl(X^{\top}\widehat MX(\Omega^2\Phi_q+\Phi_q\Omega^2)\bigr)
 +\tr(X^{\top}\widehat HX\,\Omega^2\Phi_q\Omega^2),
\]
which is the full compliance-weighted FE residual, not a residual first
projected into the trial space. Its stationarity in $Z$ is the coupled
linear system of Theorem~\ref{thm:completion}(ii),
\begin{equation}
 \widehat K_{NN}Z\Phi_q-\widehat M_{NN}Z(\Omega^2\Phi_q+\Phi_q\Omega^2)
 +\widehat H_{NN}Z\,\Omega^2\Phi_q\Omega^2=-C,
 \label{eq:coefficient-solve}
\end{equation}
where the subscript $NN$ marks the hidden block and $C$ collects the same
three terms with the port block $[E]W$. Conjugate gradients solve it with
a block-diagonal preconditioner: column $j$ uses the inverse of
$(\Phi_q)_{jj}(\widehat K_{NN}-2\omega_j^2\widehat M_{NN}+\omega_j^4\widehat H_{NN})$,
the part of the operator that does not couple modes. With
$\Phi_q=\Omega^{-2}$ the system decouples and reproduces the earlier
per-mode undamped solution to $5\times10^{-12}$. For the shell the
equality-only solve reaches a relative residual below $10^{-11}$, and the
full-residual and projected values of $J_F$ agree to
\LabProjectionAgreement.

A thin QR factorization $R[E,N]=U_D[V_E,V_N]$ with $U_D^{\top}U_D=I$ gives
$X(Z)=V_EW+V_NZ$ with the singular values of $RT$. Introduce an auxiliary
$Y=X(Z)$ with $\sigma_r(Y)\ge\delta_{\rm g}=1.001\delta$ and a scaled
multiplier $U$. The alternating-direction updates \citep{boyd2011} are
\begin{align}
 \mathcal Q(Z^{k+1})+\rho V_N^{\top}V_NZ^{k+1}
   &=-C+\rho V_N^{\top}(Y^k-U^k-V_EW),\nonumber\\
 Y^{k+1}&=\mathcal P_{\delta_{\rm g}}(X(Z^{k+1})+U^k),\qquad
 U^{k+1}=U^k+X(Z^{k+1})-Y^{k+1},
 \label{eq:admm}
\end{align}
where $\mathcal Q$ is the left side of Eq.~\eqref{eq:coefficient-solve}
and $\mathcal P_{\delta_{\rm g}}$ raises every singular value below
$\delta_{\rm g}$ to $\delta_{\rm g}$. The $Z$ step is the same
preconditioned system, warm started. The physical iterates remain
$T=EW+NZ$, so $AT=W$ holds throughout, and only an actual $X(Z)$ with
$\sigma_r\ge\delta$ is kept as a feasible candidate.

The start is $T_{\rm LS}$ with the singular values of $X(Z_{\rm LS})$
raised to $1.02\delta$ and mapped back to $Z$ by least squares. The penalty
starts at the median of the diagonal blocks' traces divided by
$\tr(V_N^{\top}V_N)$; every 25 iterations it doubles if the relative primal
residual $\|X-Y\|_F/\max(\|X\|_F,1)$ exceeds $10^{-5}$ and halves if the
scaled dual residual exceeds ten times the primal one. Both residuals have
tolerance $2\times10^{-7}$ and the budget is 1200 iterations. The
constraint is nonconvex, so convergence of the residuals does not prove
global optimality. No stopping decision uses excluded responses.

An accepted map must pass four numerical tests: aggregate and per-column
port mismatch below $10^{-10}$ relative, $\|VT-I_r\|_F<10^{-9}$, actual
$\sigma_r(RT)\ge\delta$, and numerical rank $r$ against
$10^{-10}\sigma_1$. The guards $\kappa_2(RT)<10^5$ and $\kappa_2(P)<10^6$
are inactive in both examples.

\subsection{Comparison with the undamped criterion}\label{sec:completion-benchmark}
The undamped criterion of free vibration replaces $\Phi_q$ by $\Omega^{-2}$.
Completed with it in the same trial space and under the same margin, the
shell map has $J_F=\LabJFOfJZero$ against \LabJF\ for the accepted map, and
the two maps differ by \LabMapChange\% in relative Frobenius norm. Its
excluded-response error is \LabExcludedErrorJZero\% against
\LabExcludedError\%. For the chain the two maps differ by \SpringMapChange\%.

\subsection{Coordinate selection and state-space outputs}\label{sec:state-details}
For the square form of Theorem~\ref{thm:realization}(iii), the rows of $\Pi$
are selected by column-pivoted QR of $(RT_\theta)^{\top}$ followed by
determinant-increasing row exchanges with expansion-coefficient threshold
1.05 \citep{golub2013}, with $B_\theta=L^{\top}B$ and $S_\theta=SL$. With state
$\xi=[x^{\top},\dot x^{\top}]^{\top}\in\R^{2r}$,
\begin{equation}
 \dot\xi=\underbrace{\begin{bmatrix}0&I_r\\-M_\theta^{-1}K_\theta&-M_\theta^{-1}C_\theta\end{bmatrix}}_{A_{\rm s}}\xi
 +\underbrace{\begin{bmatrix}0\\M_\theta^{-1}B_\theta\end{bmatrix}}_{B_{\rm s}}f,
 \label{eq:state}
\end{equation}
with displacement and velocity outputs $[S_\theta,0]\xi$ and $[0,S_\theta]\xi$.
Substituting the physical equation into the acceleration gives
\begin{equation}
 a=-S_\theta[\,M_\theta^{-1}K_\theta,\ M_\theta^{-1}C_\theta\,]\xi
   +S_\theta M_\theta^{-1}B_\theta f,\qquad
 S_\theta M_\theta^{-1}B_\theta=R_{\rm o}WW^{\top}R_{\rm i}^{\top}
 =\lim_{|s|\to\infty}H_\theta(s).
 \label{eq:feedthrough}
\end{equation}
The acceleration feedthrough is essential and is checked separately.
Numerical implementations use solves or congruent products rather than an
explicit $M_\theta^{-1}$. A new load column $B'$ or output row $S'$ enters
through $L^{\top}B'$ or $S'L$, which reproduces Theorem~\ref{thm:realization}(ii)
because $LP=T_\theta$.
